% Part: incompleteness
% Chapter: representability-in-q
% Section: comp-representable

\documentclass[../../../include/open-logic-section]{subfiles}

\begin{document}

\olfileid[fr]{inc}{req}{crq}
\olsection{Les fonctions calculables sont représentables dans $\Th{Q}$}

\begin{thm}
Toute fonction totale calculable est représentable dans~$\Th{Q}$.
\end{thm}

\begin{proof}
Pour fixer les idées, et conformément à la thèse de
Church--Turing, appelons fonctions totales calculables les fonctions
récursives générales. Celles-ci sont les fonctions que l'on peut obtenir à
partir de la fonction nulle~$\Zero$, du successeur~$\Succ$ et des
projections~$\Proj{n}{i}$ par composition, récurrence primitive
et minimisation régulière. D'après le
\olref[pri]{lem:prim-rec}, toute fonction~$h$ obtenue par
récurrence primitive à partir de $f$ et de~$g$ peut également
être définie, à partir de $f$, $g$, $\Zero$, $\Succ$,
$\Proj{n}{i}$, $\Add$, $\Mult$ et $\Char{=}$, au moyen de
la composition et de la minimisation régulière. Une fonction
est donc récursive générale si et seulement si elle peut être
obtenue à partir de $\Zero$, $\Succ$, $\Proj{n}{i}$, $\Add$,
$\Mult$ et $\Char{=}$ par ces deux opérations.

Nous avons montré que chacune de ces fonctions de base est
représentable dans~$\Th{Q}$
(\cref{inc:req:bre:prop:rep-zero,inc:req:bre:prop:rep-succ,inc:req:bre:prop:rep-proj,inc:req:bre:prop:rep-id,inc:req:bre:prop:rep-add,inc:req:bre:prop:rep-mult}).
Nous avons aussi montré que la composition et la minimisation
régulière de fonctions représentables restent représentables
(\olref[cmp]{prop:rep-composition},
\olref[min]{prop:rep-minimization}). Toute fonction récursive
générale est donc représentable dans~$\Th{Q}$.
\end{proof}

\begin{explain}
Nous avons ainsi caractérisé les fonctions calculables par leur
représentabilité dans~$\Th{Q}$. La portée de l'argument est en
réalité plus large. D'après la définition de la représentabilité,
toute théorie qui étend~$\Th{Q}$, ou dans laquelle on peut
interpréter~$\Th{Q}$, peut représenter les fonctions calculables.
Réciproquement, supposons qu'une théorie cohérente démontre
l'inégalité de deux numéraux distincts et que ses preuves
soient effectivement vérifiables. Si une formule $!A_f$ y
représente une fonction totale~$f$, énumérons les paires
constituées d'un entier naturel $m$ et d'un code de preuve.
Pour les arguments fixés, cherchons dans cette énumération une
preuve de l'instance de~$!A_f$ de sortie~$\num m$. Il existe
une telle preuve pour la valeur de~$f$, d'après la première
condition de représentabilité. Si une preuve donnait une
autre sortie, la seconde condition permettrait de démontrer
l'égalité de deux numéraux distincts, en contradiction avec la
cohérence. Le premier témoin rencontré donne donc la valeur
de~$f$ et la recherche s'arrête pour tous les arguments.
On peut donc, par exemple,
caractériser les fonctions calculables comme celles qui sont
représentables dans l'arithmétique de Peano, ou même dans la
théorie des ensembles de Zermelo--Fraenkel, pourvu que ces théories soient cohérentes.

Comme Gödel l'a remarqué, ce résultat peut surprendre. Nous
verrons que, pour la démontrabilité, le choix de la théorie est
déterminant : en gros, des axiomes plus forts permettent de
démontrer davantage d'énoncés. Pourtant, dans une large classe
de théories axiomatiques, les fonctions représentables sont
exactement les fonctions calculables. Renforcer les axiomes
n'ajoute aucune fonction représentable tant que la théorie
reste cohérente, distingue les numéraux et dispose d'un système
de démonstration effectif.
\end{explain}

Une théorie incohérente fournit un cas limite : elle démontre
toutes les instances exigées par les deux conditions de
représentabilité, y compris pour une fonction non calculable.

\end{document}
